\documentclass[letterpaper]{article}
\usepackage{multicol}
\usepackage[left=2cm,top=2cm,bottom=2cm,right=2cm,nohead,nofoot]{geometry}
\usepackage{enumerate}
\renewcommand{\baselinestretch}{4}
\begin{document}

1TI      3:16  \LARGE $\kappa\alpha\iota$ \normalsize  2532 {CONJ}  \LARGE $o\mu o\lambda o\gamma o\upsilon\mu\epsilon\nu\omega\varsigma$ \normalsize  3672 {ADV}  \LARGE $\mu\epsilon\gamma\alpha$ \normalsize  3173
{A-NSN}  \LARGE $\epsilon\sigma\tau\iota\nu$ \normalsize  1510 5719 {V-PAI-3S}  \LARGE $\tau o$ \normalsize  3588 {T-NSN}  \LARGE $\tau\eta\varsigma$ \normalsize  3588
{T-GSF}  \LARGE $\epsilon\upsilon\sigma\epsilon\beta\epsilon\iota\alpha\varsigma$ \normalsize  2150 {N-GSF}  \LARGE $\mu\upsilon\sigma\tau\eta\rho\iota o\nu$ \normalsize  3466 {N-NSN}  \LARGE $\theta\epsilon o\varsigma$ \normalsize  2316
{N-NSM}  \LARGE $\epsilon\phi\alpha\nu\epsilon\rho\omega\theta\eta$ \normalsize  5319 5681 {V-API-3S}  \LARGE $\epsilon\nu$ \normalsize  1722 {PREP}  \LARGE $\sigma\alpha\rho\kappa\iota$ \normalsize  4561
{N-DSF}  \LARGE $\epsilon\delta\iota\kappa\alpha\iota\omega\theta\eta$ \normalsize  1344 5681 {V-API-3S}  \LARGE $\epsilon\nu$ \normalsize  1722 {PREP}  \LARGE $\pi\nu\epsilon\upsilon\mu\alpha\tau\iota$ \normalsize 
4151 {N-DSN}  \LARGE $\omega\phi\theta\eta$ \normalsize  3708 5681 {V-API-3S}  \LARGE $\alpha\gamma\gamma\epsilon\lambda o\iota\varsigma$ \normalsize  32 {N-DPM}
 \LARGE $\epsilon\kappa\eta\rho\upsilon\chi\theta\eta$ \normalsize  2784 5681 {V-API-3S}  \LARGE $\epsilon\nu$ \normalsize  1722 {PREP}  \LARGE $\epsilon\theta\nu\epsilon\sigma\iota\nu$ \normalsize  1484 {N-DPN}
 \LARGE $\epsilon\pi\iota\sigma\tau\epsilon\upsilon\theta\eta$ \normalsize  4100 5681 {V-API-3S}  \LARGE $\epsilon\nu$ \normalsize  1722 {PREP}  \LARGE $\kappa o\sigma\mu\omega$ \normalsize  2889 {N-DSM}
 \LARGE $\alpha\nu\epsilon\lambda\eta\phi\theta\eta$ \normalsize  353 5681 {V-API-3S}  \LARGE $\epsilon\nu$ \normalsize  1722 {PREP}  \LARGE $\delta o\xi\eta$ \normalsize  1391 {N-DSF}
     \renewcommand{\baselinestretch}{1}
\begin{multicols}{2}
\small
\section*{Vocabulary}
\begin{description}

\item[32: ]



1) a messenger, envoy, one who is sent, an angel, a messenger
from God



\item[353: ]



1) to take up, raise

2) to take up (a thing in order to carry or use it)



\item[1344: ]



1) to render righteous or such he ought to be

2) to show, exhibit, evince, one to be righteous, such as he is
and wishes himself to be considered

3) to declare, pronounce, one to be just, righteous, or such as
he ought to be



\item[1391: ]



1) opinion, judgment, view

2) opinion, estimate, whether good or bad concerning someone

2a) in the NT always a good opinion concerning one, resulting in praise, honour, and glory

3) splendour, brightness

4) a most glorious condition, most exalted state

4a) of that condition with God the Father in heaven to which
Christ was raised after he had achieved his work on earth

4b) the glorious condition of blessedness into which is
appointed and promised that true Christians shall enter
after their Saviour's return from heaven



\item[1484: ]



1) a multitude (whether of men or of beasts) associated or
living together

2) a multitude of individuals of the same nature or genus

3) a tribe, nation, people group

4) in the OT, foreign nations not worshipping the true God, pagans,
Gentiles

5) Paul uses the term for Gentile Christians
For Synonyms see entry 5927



\item[1510: ]



1) to be, to exist, to happen, to be present



\item[1722: ]



1) in, by, with etc.



\item[2150: ]



1) reverence, respect

2) piety towards God, godliness



\item[2316: ]


1) a god or goddess, a general name of deities or divinities

2) the Godhead, trinity

2a) God the Father, the first person in the trinity

2b) Christ, the second person of the trinity

2c) Holy Spirit, the third person in the trinity

3) spoken of the only and true God

4) whatever can in any respect be likened unto God, or resemble him in
any way


\item[2532: ]



1) and, also, even, indeed, but



\item[2784: ]



1) to be a herald, to officiate as a herald

2) to publish, proclaim openly: something which has been done

3) used of the public proclamation of the gospel and matters
pertaining to it, made by John the Baptist, by Jesus, by the
apostles and other Christian teachers



\item[2889: ]



1) an apt and harmonious arrangement or constitution, order,
government

2) ornament, decoration, adornment, i.e. the arrangement of the stars,
'the heavenly hosts' , as the ornament of the heavens. 1Pe 3:3

3) the world, the universe

4) the circle of the earth, the earth

5) the inhabitants of the earth, men, the human family

6) the ungodly multitude; the whole mass of men alienated from God,

and therefore hostile to the cause of Christ

7) world affairs, the aggregate of things earthly

8) any aggregate or general collection of particulars of any sort



\item[3173: ]



1) great

1a) of the external form or sensible appearance of things (or
of persons)

1b) of number and quantity: numerous, large, abundant

1c) of age: the elder

1d) used of intensity and its degrees: with great effort, of
the affections and emotions of the mind, of natural events
powerfully affecting the senses: violent, mighty, strong

2) predicated of rank, as belonging to

2a) persons, eminent for ability, virtue, authority, power

2b) things esteemed highly for their importance: of great
moment, of great weight, importance

2c) a thing to be highly esteemed for its excellence: excellent

3) splendid, prepared on a grand scale, stately

4) great things



\item[3466: ]



1) hidden thing, secret, mystery

1a) generally mysteries, religious secrets, confided only to the
initiated and not to ordinary mortals

1b) a hidden or secret thing, not obvious to the understanding

1c) a hidden purpose or counsel

2) in rabbinic writings, it denotes the mystic or hidden sense

2a) of an OT saying

2b) of an image or form seen in a vision

2c) of a dream



\item[3588: ]



1) this, that, these, etc.
Only significant renderings other than "the" counted



\item[3672: ]



1) by consent of all, confessedly, without controversy



\item[3708: ]



1) to see with the eyes

2) to see with the mind, to perceive, know

3) to see, i.e. become acquainted with by experience, to experience

4) to see, to look to

5) I was seen, showed myself, appeared.
For Synonyms see entry 5822



\item[4100: ]



1) to think to be true, to be persuaded of, to credit, place
confidence in

1a) of the thing believed

1b) in a moral or religious reference

2) to entrust a thing to one, i.e. his fidelity


\item[4151: ]



1) the third person of the triune God, the Holy Spirit, coequal,
coeternal with the Father and the Son


2) the spirit, i.e. the vital principal by which the body is animated

3) a spirit, i.e. a simple essence, devoid of all or at least
all grosser matter, and possessed of the power of knowing,
desiring, deciding, and acting

4) the disposition or influence which fills and governs the soul
of any one

5) a movement of air (a gentle blast)



\item[4561: ]



1) flesh (the soft substance of the living body, which covers
the bones and is permeated with blood) of both man and beasts

2) the body

2a) the body of a man

2b) used of natural or physical origin, generation or
relationship

2c) the sensuous nature of man, "the animal nature"

3) a living creature (because possessed of a body of flesh)
whether man or beast

4) the flesh, denotes mere human nature, the earthly nature of man apart
from divine influence, and therefore prone to sin and opposed to God



\item[5319: ]



1) to make manifest or visible or known what has been hidden or unknown,
to manifest, whether by words, or deeds, or in any other way

1a) make actual and visible, realised

1b) to make known by teaching

1c) to become manifest, be made known

1d) of a person

1e) to become known, to be plainly recognised, thoroughly understood



\end{description}
\section*{Grammar}
\begin{description}

\item[5681]
Aorist Tense, Passive  Voice, Indicative  Mood
\item[5719]
Present Tense, Indicative  Mood
\item[Aorist Tense:]
The aorist tense is characterized by its emphasis on punctiliar action; that is, the concept of the verb is considered without regard for past, present, or future time.  There is no direct or clear English equivalent for this tense, though it is generally rendered as a simple past tense in most translations.

\item[Indicative  Mood:]
The indicative mood is a simple statement of fact.  If an action really occurs or has occurred or will occur, it will be rendered in the indicative mood.

\item[Passive  Voice:]
The passive voice represents the subject as being the recipient of the action.  E.g., in the sentence, "The boy was hit by the ball," the boy receives the action.

\item[Present Tense:]
The present tense represents a simple statement of fact or reality viewed as occurring in actual time.  In most cases this corresponds directly with the English present tense.

\end{description}
\end{multicols}

\begin{verbatim}
KEY TO PARSINGS:

  UNDECLINED FORMS:
      ADV   = ADVerb or adverb and particle combined
      CONJ  = CONJunction or conjunctive particle
      COND  = CONDitional particle or conjunction
      PRT   = PaRTicle, disjunctive particle
      PREP  = PREPosition

  DECLINED FORMS: prefix-case-number-gender-(suffix)
    Prefixes:
      N-    = Noun
      A-    = Adjective
      T-    = definite arTicle

    Cases (5-case system only):
      -N    = Nominative
      -V    = Vocative
      -G    = Genitive
      -D    = Dative
      -A    = Accusative

    Number:                       Gender:
      S     = Singular              M     = Masculine
      P     = Plural                F     = Feminine

  VERBS: V-tense-voice-mood-person-number
      V-PAI-3S = Present Active Indicative, 3rd Person Singular
      V-API-3S = Aorist Passive Indicative, 3rd Person Singular


Original raw source:
     3:16 kai 2532 {CONJ} omologoumenwv 3672 {ADV} mega 3173
{A-NSN} estin 1510 5719 {V-PAI-3S} to 3588 {T-NSN} thv 3588
{T-GSF} eusebeiav 2150 {N-GSF} musthrion 3466 {N-NSN} yeov 2316
{N-NSM} efanerwyh 5319 5681 {V-API-3S} en 1722 {PREP} sarki 4561
{N-DSF} edikaiwyh 1344 5681 {V-API-3S} en 1722 {PREP} pneumati
4151 {N-DSN} wfyh 3708 5681 {V-API-3S} aggeloiv 32 {N-DPM}
ekhrucyh 2784 5681 {V-API-3S} en 1722 {PREP} eynesin 1484 {N-DPN}
episteuyh 4100 5681 {V-API-3S} en 1722 {PREP} kosmw 2889 {N-DSM}
anelhfyh 353 5681 {V-API-3S} en 1722 {PREP} doxh 1391 {N-DSF}
\end{verbatim}


\end{document}
